\medskip

On considère le cube ABCDEFGH.
On note I le milieu du segment[EH] et on considère le triangle CFI.
L'espace est muni du repère orthonormé $\left(\text{A}~;~ \vect{\text{AB}},\: \vect{\text{AD}},\: \vect{\text{AE}}\right)$ et on admet que le point I a pour coordonnées $\left(0~;~\dfrac12;~1\right)$ dans ce repère.

\begin{center}
\psset{unit=1cm}
\begin{pspicture}(7.4,8)
\pspolygon(0.2,0.8)(4.6,0.2)(4.6,4.6)(0.2,5.2)
\psline(4.6,0.2)(7.2,2.6)(7.2,7)(4.6,4.6)
\psline(7.2,7)(2.8,7.6)(0.2,5.2)
\psline[linestyle=dashed](0.2,0.8)(2.8,3.2)(7.2,2.6)
\psline[linestyle=dashed](2.8,3.2)(2.8,7.6)
\psline(5,7.3)(0.2,5.2)(4.6,0.2)
\psline[linestyle=dashed](5,7.3)(4.6,0.2)
\psdots(2.8,3.2)(0.2,0.8)(4.6,0.2)(7.2,2.6)(2.8,7.6)(0.2,5.2)(4.6,4.6)(7.2,7)(5,7.3)(2.4,4.9)
\uput[l](2.8,3.2){A} \uput[l](0.2,0.8){B} \uput[dr](4.6,0.2){C} \uput[r](7.2,2.6){D}
\uput[u](2.8,7.6){E} \uput[l](0.2,5.2){F} \uput[dl](4.6,4.6){G} \uput[ur](7.2,7){H}
\uput[u](5,7.3){I} \uput[u](2.4,4.9){J}
\end{pspicture}
\end{center}

\smallskip

\begin{enumerate}
\item 
	\begin{enumerate}
		\item Les coordonnées des points C, F et G sont:
$\text{C}\begin{pmatrix}1\\1\\0\end{pmatrix}$, $\text{F}\begin{pmatrix}1\\0\\1\end{pmatrix}$ et $\text{G}\begin{pmatrix}1\\1\\1\end{pmatrix}$
		
		\item Soit le vecteur $\vect{n}\begin{pmatrix}1\\2\\2\end{pmatrix}$.
		
\begin{list}{\textbullet}{}
\item $\vectt{CF}\begin{pmatrix}0\\-1\\1\end{pmatrix}$;
$\vect{n}\cdot \vectt{CF} = 1\times 0+2\times (-1) + 2\times 1=0$ donc $\vect{n}\perp \vectt{CF}$
\item $\vectt{CI}\begin{pmatrix}-1\\-\frac{1}{2}\\1\end{pmatrix}$;
$\vect{n}\cdot \vectt{CI} = 1\times (-1)+2\times (-\frac{1}{2}) + 2\times 1=0$ donc $\vect{n}\perp \vectt{CI}$
\item Les points C, F et I ne sont pas alignés donc ces trois points forment le plan (CFI) dont $\vectt{CF}$ et $\vectt{CI}$ sont deux vecteurs directeurs.
\end{list}

Le vecteur $\vect{n}$ est orthogonal à deux vecteurs directeurs du plan (CFI) donc il est normal au plan (CFI).
	
		\item
D'après la question précédente, le plan (CFI) a une équation de la forme \\
$x+2y+2z+d=0$ où $d\in\R$.

$\text{C} \in (\text{CFI})$ donc $x_{\text{C}}+2y_{\text{C}}+2z_{\text{C}}+d=0$ soit $1+2\times 1 +2\times 0+d=0$ donc $d=-3$.

Une équation cartésienne du plan (CFI) est donc: $x + 2y + 2z - 3 = 0$.		

	\end{enumerate}	
\item  On note $d$ la droite passant par G et orthogonale au plan (CFI).
	\begin{enumerate}
		\item 
La droite $d$ est orthogonale au plan (CFI); elle a donc le vecteur $\vect{n}$ comme vecteur directeur.
De plus elle passe par le point G.
C'est donc l'ensemble des points M de coordonnées $(x\;; y\;; z)$ tels que $\vectt{GM}$ et $\vect n$ soient colinéaires, c'est-à-dire tels que $\vectt{GM} =t.\vect n$ avec $t\in\R$.
		
$\vectt{GM} =t.\vect n
\iff
\left \{
\begin{array}{l !{=} r}
x-1 & t\\
y-1 & 2t\\
z-1 & 2t
\end{array}
\right . $		

La droite $d$ a donc pour représentation paramétrique:
$\left \{
\begin{array}{l !{=} r}
x & 1+t\\
y & 1+2t\\
z &1+ 2t
\end{array}
\right . \text{ avec }t\in\R$
		
		\item 
La droite $d$ est orthogonale au plan (CFI) donc le projeté orthogonal K du point G de la droite $d$ sur le plan (CFI), est le point d'intersection de $d$ et de (CFI); donc ses coordonnées $(x\;; y\;; z)$ vérifient le système:
$\left \{
\begin{array}{r !{=} l}
x & 1+t\\
y & 1+2t\\
z &1+ 2t\\
x+2y+2z-3 & 0
\end{array}
\right .$
		
On a donc: $(1+t) +2\left (1+2t\right ) + 2\left (1+2t\right ) -3=0$,
ce qui donne: \\
$1+t+2+4t+2+4t-3=0$
soit
$9t=-2$ donc $t=-\frac{2}{9}$.

$\left \{
\begin{array}{r !{=} r !{=} l !{=} l}
x & 1+t & 1 -\frac{2}{9} & \frac{7}{9}\\[4pt]
y & 1+2t & 1- 2\times \frac{2}{9} & \frac{5}{9}\\[4pt]
z &1+ 2t & 1- 2\times \frac{2}{9} & \frac{5}{9}
\end{array}
\right .$

Donc le projeté orthogonal K du point G sur le plan (CFI) a pour coordonnées $\left (\dfrac{7}{9}\;; \dfrac{5}{9}\;; \dfrac{5}{9} \right )$.
		\item La distance du point G au plan (CFI) est:
		
$\begin{gathered}\text{GK}= \ds\sqrt{(x_{\text K} - x_{\text G})^2 + (y_{\text K} - y_{\text G})^2 +(z_{\text K} - z_{\text G})^2 }
= \ds\sqrt{\left ( \frac{7}{9} - 1\right )^2 + \left ( \frac{5}{9} - 1\right )^2+\left ( \frac{5}{9} - 1\right )^2}\\
\phantom{\text{GK}}
= \ds\sqrt{\left ( -\frac{2}{9}\right )^2 + \left ( -\frac{4}{9}\right )^2+\left ( -\frac{4}{9}\right )^2}
=\ds\sqrt{\frac{4}{81} + \frac{16}{81} + \frac{16}{81}}
=\ds\sqrt{ \frac{36}{81}}
= \frac{6}{9}
= \frac23\end{gathered}$.
	\end{enumerate}
	
\item On considère la pyramide GCFI.







	\begin{enumerate}
		\item 
Si on considère que I est le sommet de la pyramide, la base en est le triangle GFC rectangle en G.

Le projeté orthogonal de I sur le plan (GFC) est le point J, milieu de [FG], et $\text{IJ}=1$ donc la hauteur $h$ vaut 1.

Le triangle GFC a pour aire $\dfrac{\text{GF}\times \text{GC}}{2} = \dfrac{1}{2}$.

La pyramide GCFI a donc pour volume, en unité de volume:  $\dfrac{1}{3}\times \dfrac{1}{2}\times 1 = \dfrac{1}{6}$.
		
		\item Appelons $\mathcal{A}$ l'aire du triangle CFI.
		
Si on considère que G est le sommet de la pyramide, la base en est le triangle CFI d'aire $\mathcal{A}$.

Le projeté orthogonal de G sur le plan (CFI) est le point K, donc la hauteur $h$ vaut GK soit $\dfrac{2}{3}$.

La pyramide GCFI a donc pour volume:  $\dfrac{1}{3}\times \mathcal{A}\times \dfrac{2}{3} = \dfrac{2\mathcal{A}}{9}$.

Or le volume de la pyramide vaut $\dfrac{1}{6}$, donc $\dfrac{2\mathcal{A}}{9} = \dfrac{1}{6}$ donc $\mathcal{A}=\dfrac{9}{12}=\dfrac{3}{4}$.

L'aire du triangle CFI est, en unité d'aire, égale à $\dfrac{3}{4}$.
		
	\end{enumerate}
\end{enumerate}
